\chapter{Parser}
\section{Lexing: Phân tích Từ vựng chuyên sâu (Lexer)}

`Lexer.cs` là thành phần đầu tiên của máy chủ KBMS tiếp nhận truy vấn người dùng. Nó thực hiện quét chuỗi (Scanning) $O(L)$ để bẻ gãy văn bản thô thành một danh sách các nguyên tử ngôn ngữ gọi là \textbf{Tokens}.

\subsection{Cơ chế Quét (Scanning Mechanism)}

Lexer vận hành như một máy trạng thái (State Machine) với hai con trỏ: `\_start` (điểm bắt đầu của một từ) và `\_current` (vị trí đang quét hiện tại).

![lexer\_token\_flow\_v2.png](../assets/diagrams/lexer\_token\_flow\_v2.png)
\textit{Hình: Quy trình quét chuỗi và định danh Token (dàn ngang)}

\subsubsection{Quy trình xử lý (`ScanToken`):}
1.  \textbf{Bỏ qua nhiễu}: Loại bỏ khoảng trắng (` `), tab (`\t`), xuống dòng (`\n`). Khi gặp xuống dòng, Lexer tự động tăng chỉ số `\_line` và reset `\_column` để phục vụ báo lỗi.
2.  \textbf{Ký tự đơn}: Định danh ngay lập tức các Token như `(`, `)`, `,`, `;`, `+`, `-`, `*`, `/`.
3.  \textbf{Toán tử kép}: Sử dụng hàm `Match()` để kiểm tra ký tự tiếp theo (Ví dụ: `>` và `>=`).
4.  \textbf{Chú thích (Comments)}: Khi gặp `--`, Lexer sẽ bỏ qua toàn bộ phần còn lại của dòng.

---

\subsection{Xử lý các Hằng số (Literals)}

`Lexer.cs` chứa các phương thức chuyên biệt để bóc tách dữ liệu thực tế:

\subsubsection{Hằng số Số (`NumberLiteral`)}
Hỗ trợ đầy đủ các định dạng số học:
*   \textbf{Số nguyên}: `123`.
*   \textbf{Số thực}: `123.45`.
*   \textbf{Số âm}: Nhận diện dấu `-` đứng trước chữ số.
*   \textbf{Số mũ (Scientific Notation)}: Hỗ trợ `1.23e-10`.
*   \textbf{Xử lý nội bộ}: Sử dụng `double.TryParse` với `CultureInfo.InvariantCulture` để đảm bảo tính nhất quán trên mọi hệ điều hành.

\subsubsection{Hằng số Chuỗi (`StringLiteral`)}
Hỗ trợ cả nháy đơn `'` và nháy kép `"`:
*   \textbf{Ký tự thoát (Escape)}: Cho phép dùng nháy kép lặp lại (`""`) để biểu diễn một dấu nháy đơn trong chuỗi.
*   \textbf{Đa dòng}: Chuỗi có thể kéo dài qua nhiều dòng, Lexer sẽ tự nạp các ký tự xuống dòng vào giá trị literal.

---

\subsection{Quản lý Từ khóa (Keyword Dictionary)}

KBMS sở hữu một kho từ điển từ khóa (Keywords) đồ sộ với hơn \textbf{190 Token Types}.

| Nhóm | Ví dụ từ khóa |
| :--- | :--- |
| \textbf{DML} | `SELECT`, `INSERT`, `UPDATE`, `DELETE`, `SOLVE` |
| \textbf{DDL} | `CREATE`, `DROP`, `ALTER`, `CONCEPT`, `RELATION` |
| \textbf{Types} | `INT`, `DECIMAL`, `VARCHAR`, `BOOLEAN`, `DATE` |
| \textbf{TCL} | `BEGIN`, `COMMIT`, `ROLLBACK` |

Lexer thực hiện tra cứu không phân biệt hoa thường (`OrdinalIgnoreCase`) để tạo ra các Token có Type tương ứng, giúp Parser xử lý nhanh chóng mà không cần so khớp chuỗi (String comparison).

> [!CAUTION]
> \textbf{ParseException trong Lexer}
> Nếu Lexer gặp một ký tự không mong muốn (ví dụ: `!`), nó sẽ ném ra một `ParseException` kèm tọa độ `Line/Column`. Điều này giúp hệ thống báo lỗi ngay từ tầng từ vựng trước khi chuyển sang tầng cú pháp. 🚨

\section{Tầng Parser: Cơ chế Đệ quy Đi xuống (Recursive Descent)}

Lớp `Parser.cs` là "ngã tư" điều hướng mọi câu lệnh. Với hơn 3400 dòng mã nguồn, nó thực hiện một quy trình phân tích cú pháp hiệu năng cao dựa trên thuật toán \textbf{Recursive Descent (Đệ quy đi xuống)}.

\subsection{Cơ chế Điều hướng (Sentence Dispatching)}

Mỗi khi máy chủ nhận được một chuỗi Tokens từ Lexer, nó khởi tạo `Parser.cs` và gọi hàm `ParseAll()`.

\begin{verbatim}
// Luồng xử lý đa câu lệnh trong Parser.cs
public List<AstNode> ParseAll() {
    var statements = new List<AstNode>();
    while (!IsAtEnd()) {
        var stmt = ParseStatement(); // "Ngã tư" điều hướng
        statements.Add(stmt);
        if (Check(TokenType.SEMICOLON)) Advance(); // Tiêu thụ dấu ';'
    }
}
\end{verbatim}

\subsubsection{Hàm `ParseStatement()`}
Đây là bộ khung switch-case khổng lồ, quyết định câu lệnh tiếp theo là DDL (Định nghĩa), DML (Truy vấn) hay TCL (Giao dịch). Mỗi khi nhận diện được một từ khóa bắt đầu (như `CREATE`, `SELECT`), nó sẽ ủy quyền cho một hàm con chuyên biệt (ví dụ: `ParseCreate()`, `ParseSelect()`).

---

\subsection{Kỹ thuật Dự đoán & Tiêu thụ (Look-ahead)}

KBMS sử dụng ngữ pháp \textbf{LL(k)}, cho phép Parser "nhìn trước" $k$ Token để quyết định nhánh thực thi mà không cần quay lui (Backtracking).

![ast\_tree\_layout\_v2.png](../assets/diagrams/ast\_tree\_layout\_v2.png)
\textit{Hình: Ví dụ cấu trúc cây AST sau khi phân tích (dàn ngang)}

*   \textbf{`Peek()`}: Xem Token hiện tại mà không tiêu thụ.
*   \textbf{`Check(type)`}: Kiểm tra xem Token hiện tại có đúng Type mong muốn hay không.
*   \textbf{`Consume(type)`}: Nếu Token đúng Type, tiến tới Token tiếp theo. Nếu sai, ngay lập tức ném ra `ParseException` kèm tọa độ `Line/Column`.
*   \textbf{`Advance()`}: Luôn tiến về phía trước, đảm bảo độ phức tạp thời gian là $O(N)$ (N là số lượng Token).

---

\subsection{Phân tích Khối Lệnh phức tạp (Complex Parsing)}

Một trong những phần xử lý tinh vi nhất trong `Parser.cs` là phân tích lệnh `CREATE CONCEPT`. Thay vì một parser phẳng, nó thực hiện đệ quy lồng nhau:

1.  \textbf{`ParseCreate`} gọi \textbf{`ParseCreateConcept`}.
2.  \textbf{`ParseCreateConcept`} lặp trong khối `(...)` và kiểm tra từ khóa.
3.  Nếu gặp `VARIABLES`, nó gọi \textbf{`ParseVariableDefinition`} để bóc tách `name: TYPE`.
4.  Nếu gặp `CONSTRAINTS`, nó gọi \textbf{`ParseConstraintList`} (đây là lúc Expression Parser được kích hoạt).
5.  Nếu gặp `RULES`, nó gọi \textbf{`ParseConceptRuleList`} (phân tích IF-THEN logic).

> [!CAUTION]
> \textbf{Error Reporting \& Recovery}
> Khi gặp một lỗi cú pháp (ví dụ: thiếu dấu ngoặc đóng), Parser sẽ ném một `ParseException`. Hệ thống thu nạp `Line` và `Column` từ chính Token bị lỗi. Điều này cực kỳ quan trọng giúp nhà phát triển tìm ra và sửa lỗi ngay lập tức trên bản đồ tri thức lớn. 🕵️‍♂️

\section{Hệ thống Cây cú pháp Trừu tượng (AST Nodes)}

Cây AST là đại diện logic hoàn chỉnh cho tri thức sau khi được biên dịch. Tại KBMS, mỗi câu lệnh được bóc tách thành một hoặc nhiều `AstNode` chuyên biệt, phân cấp theo 5 họ ngôn ngữ thực tế trong thư mục `KBMS.Parser/Ast`.

![ast\_overview\_v2.png](../assets/diagrams/ast\_overview\_v2.png)
\textit{Hình: Tổng quan 5 nhánh Phả hệ Tri thức KBMS}

---

\subsection{Hệ phả KDL (Knowledge Definition Language)}
Đây là họ lớn nhất, chịu trách nhiệm định nghĩa "cơ thể" của tri thức:

![ast\_kdl\_detail.png](../assets/diagrams/ast\_kdl\_detail.png)
\textit{Hình: Chi tiết các Node định nghĩa tri thức (KDL)}

*   \textbf{Knowledge Base}: `CreateKnowledgeBaseNode`, `DropKbNode`, `UseKbNode`, `AlterKbNode`.
*   \textbf{Concepts}: `CreateConceptNode`, `DropConceptNode`, `AlterConceptNode`, `AddVariableNode`.
*   \textbf{Rules}: `CreateRuleNode`, `DropRuleNode`.
*   \textbf{Hierarchy}: `AddHierarchyNode`, `RemoveHierarchyNode`.
*   \textbf{Relation}: `CreateRelationNode`, `DropRelationNode`.
*   \textbf{Operator}: `CreateOperatorNode`, `DropOperatorNode`.
*   \textbf{Function}: `CreateFunctionNode`, `DropFunctionNode`.
*   \textbf{Advanced}: `AddComputationNode`, `RemoveComputationNode`, `CreateIndexNode`, `CreateTriggerNode`.

---

\subsection{Hệ phả KQL & KML (Query & Maintenance)}
Các nhánh này chịu trách nhiệm hỏi đáp và bảo trì dữ liệu tri thức:

![ast\_kql\_kml\_detail.png](../assets/diagrams/ast\_kql\_kml\_detail.png)
\textit{Hình: Chi tiết các Node Truy vấn (KQL) và Thao tác (KML)}

\subsubsection{Nhóm KQL (Knowledge Query Language)}
*   \textbf{`SelectNode`}: Dùng để lọc và truy vấn các đối tượng tri thức dựa trên thuộc tính.
*   \textbf{`SolveNode`}: Dùng cho bộ máy suy diễn (`ReasoningEngine`), chứa các mệnh đề `Given` (giả thiết) và `Find` (mục tiêu).
*   \textbf{`ShowNode`}: Liệt kê danh sách các thành phần (Concepts, Rules, KnowledgeBases).
*   \textbf{`DescribeNode`}: Xem chi tiết cấu trúc của một Concept cụ thể.
*   \textbf{`ExplainNode`}: Giải thích luồng suy luận của một câu lệnh tri thức.

\subsubsection{Nhóm KML (Knowledge Maintenance Language)}
*   \textbf{`InsertNode` / `InsertBulkNode`}: Nạp tri thức mới vào đĩa.
*   \textbf{`UpdateNode` / `DeleteNode`}: Sửa đổi và xóa bỏ các Object Instance.
*   \textbf{`ImportNode` / `ExportNode`}: Giao lưu tri thức với các tệp tin bên ngoài (`.csv`, `.json`).
*   \textbf{`MaintenanceNode`}: Các lệnh quản trị bộ nhớ thấp cấp (v.d: `REINDEX`, `VACUUM`).

---

\subsection{Hệ phả KCL & TCL (Control & Transaction)}
Nhánh bảo mật và quản lý giao dịch:

![ast\_kcl\_tcl\_detail.png](../assets/diagrams/ast\_kcl\_tcl\_detail.png)
\textit{Hình: Chi tiết các Node Bảo mật (KCL) và Giao dịch (TCL)}

*   \textbf{DCL (Control)}: `CreateUserNode`, `DropUserNode`, `AlterUserNode`, `GrantNode`, `RevokeNode`.
*   \textbf{TCL (Transaction)}: `BeginTransactionNode`, `CommitNode`, `RollbackNode`.

---

\subsection{Đặc điểm chung của mọi Node}
Mọi lớp Node đều kế thừa từ \textbf{`AstNode.cs`}, đảm bảo các tính năng cốt lõi:
1.  \textbf{Vị trí nguồn}: Lưu trữ `Line` (Dòng) và `Column` (Cột) của câu lệnh gốc.
2.  \textbf{Type định danh}: Mọi Node đều có thuộc tính `Type` (như `"CREATE\_CONCEPT"`) giúp Server Routing nhanh chóng.
3.  \textbf{ToQueryString()}: Khả năng tái tạo lại mã nguồn KBQL từ cây AST (Reverse Engineering).

> [!CAUTION]
> \textbf{Nhất quán Toàn vẹn}
> Khi Parser tạo ra một `List<AstNode>`, mọi mối quan hệ phụ thuộc giữa chúng (Ví dụ: `CreateRelationNode` tham chiếu đến `ConceptName`) đều được giữ nguyên để phục vụ quá trình kiểm tra logic ở tầng tiếp theo. 🏗️

\section{Bộ máy Đánh giá Biểu thức (Expression Engine)}

Expression Engine là "trái tim" của mọi câu lệnh tri thức. Lớp `ExpressionNode.cs` và các lớp con trong thư mục `KBMS.Parser/Ast/Expressions` chịu trách nhiệm xử lý các phép toán và logic phức tạp.

\subsection{Phân tích Biểu thức Nhị phân (Binary Expressions)}

Đây là các phép toán có hai vế, được xử lý thông qua `BinaryExpressionNode`.

\textit{   \textbf{Toán tử Toán học}: `+`, `-`, `}`, `/`, `^`, `\%`.
*   \textbf{Toán tử So sánh}: `=`, `!=`, `<`, `>`, `<=`, `>=`.
*   \textbf{Toán tử Logic}: `AND`, `OR`.

\subsubsection{Thứ tự ưu tiên (Precedence)}
Parser của KBMS thực hiện phân tích theo mức độ ưu tiên chuẩn quốc tế để đảm bảo kết quả chính xác:
1.  \textbf{Cấp 1}: Dấu ngoặc `(...)`.
2.  \textbf{Cấp 2}: Hàm gọi và Truy cập thuộc tính (`.`).
3.  \textbf{Cấp 3}: Lũy thừa `^`.
4.  \textbf{Cấp 4}: Nhân/Chia `*`, `/`.
5.  \textbf{Cấp 5}: Cộng/Trừ `+`, `-`.
6.  \textbf{Cấp 6}: So sánh `=`, `!=`, `<`, `>`.
7.  \textbf{Cấp 7}: Logic `NOT`, `AND`, `OR`.

---

\subsection{Các Đơn vị Biểu thức Cốt lõi (Core Expression Units)}

Trong cây AST, biểu thức được phân rã thành các nút lá và nút nhánh:

\subsubsection{`VariableNode` (Biến tri thức)}
Đại diện cho các thuộc tính của Concept (Ví dụ: `a`, `age`). Parser hỗ trợ truy cập thuộc tính lồng nhau thông qua dấu chấm (`p1.x`), giúp xử lý tri thức đa tầng.

\subsubsection{`LiteralNode` (Hằng số)}
Chứa giá trị thực tế như số (`100`), chuỗi (`"Active"`), hoặc Boolean (`true`/`false`). Mỗi nút lưu trữ cả giá trị thô và kiểu dữ liệu (`TokenType`).

\subsubsection{`FunctionCallNode` (Hàm gọi)}
Xử lý các hàm toán học hoặc tri thức tích hợp (v.đ: `Sqrt(x)`, `Abs(y)`). Parser bóc tách tên hàm và một danh sách các tham số (`List<ExpressionNode>`), cho phép các hàm có thể lồng nhau linh hoạt.

---

\subsection{Xử lý Biểu thức Đơn phân (Unary Expressions)}

Được đại diện bởi `UnaryExpressionNode`, dùng cho các trường hợp chỉ có một vế:
*   \textbf{Phép phủ định}: `NOT (a > B)`.
*   \textbf{Số âm}: `-x`.

> [!IMPORTANT]
> \textbf{Khả năng Tương thích Cao}
> Expression Engine của KBMS được thiết kế để kết quả đầu ra luôn là một lớp `Expression` chuẩn hóa, giúp bộ máy suy diễn (`ReasoningEngine`) có thể thực thi các biểu thức này một cách đồng nhất trên mọi loại tri thức. 🧪

\section{Xác thực Bộ biên dịch (Parser Validation)}

Parser là thành phần nhạy cảm nhất đối với đầu vào của người dùng. Mọi cú pháp sai sót đều phải được ngăn chặn ngay lập tức.

\subsection{Thống kê Độ bao phủ Testing}

Tầng Parser được bảo vệ bởi \textbf{21,000+ kịch bản test} trong `ParserTests.cs` và `LexerTests.cs`.

*   \textbf{Positive Cases}: 100\% các từ khóa và cấu trúc KDL/KQL/KML/KCL/TCL được bao phủ.
*   \textbf{Negative Cases}: Kiểm tra khả năng từ chối các câu lệnh thiếu dấu chấm phẩy, sai kiểu dữ liệu hoặc sai tên Concept.

\subsubsection{Minh chứng Mã nguồn (Parser):}
\textbf{[Missing Image: code_test_parser.png]}\\ \textit{Hình 11.1: Minh chứng mã nguồn kiểm thử độ bao phủ của bộ phân tích (Parser).}

\subsubsection{Minh chứng Kết quả (21k Tests):}
\textbf{[Missing Image: result_test_parser.png]}\\ \textit{Hình 11.2: Kết quả vượt qua toàn bộ 21,000+ kịch bản kiểm thử cú pháp.}

\subsection{Đặc tả Lỗi Cú pháp (Syntax Error Proof)}

Hệ thống báo lỗi chi tiết đến từng tọa độ (Dòng, Cột) giúp người dùng dễ dàng hiệu chỉnh:

\begin{verbatim}
-- Lỗi mô phỏng: Thiếu dấu ngoặc đóng
CREATE CONCEPT X ( VARIABLES ( id: INT );
-- [ERROR] Unexpected token ';', expected ')' at Line 1, Col 35.
\end{verbatim}

\subsubsection{Minh chứng Lỗi (Studio Editor):}
\textbf{[Missing Image: terminal_test_parser_nested.png]}\\ \textit{Hình 11.3: Chứng minh Parser xử lý thành công các biểu thức lồng nhau phức tạp.}

---

> [!IMPORTANT]
> Toàn bộ logic giải mã biểu thức (`ExpressionEngine`) trong `full\_test.kbql` được kiểm chứng từng bước thông qua các unit test nhỏ trong `LexerTests.cs`.

